% Complete reviewed Unit 4; editorial master: chapters/ch04/student.md
% Markdown SHA-256: bcaabb3fef42b5395e107d8d371a7b0865706515b05edbc46f93d8d001108748
% One-use initial assembly followed by complete content and page review.
\newcommand{\runningtitle}{Unit 4. Strong verbs I–III}
% Layout inherited from the reviewed previous course; set a running title per unit.
% This is a small hand-maintained TeX layout, not a conversion pipeline.
\RequirePackage{fix-cm}
\documentclass[11pt,twoside]{article}
\usepackage{fontspec}
\setmainfont{Cambria}[
  Renderer=HarfBuzz,
  BoldFont=cambriab.ttf,
  ItalicFont=cambriai.ttf,
  BoldItalicFont=cambriaz.ttf
]
\setmonofont{lmmono10-regular.otf}[BoldFont=lmmonolt10-bold.otf]
\usepackage[paperwidth=165mm,paperheight=241mm,inner=21mm,outer=19mm,
  top=22mm,bottom=20mm,headheight=14pt,headsep=7mm,footskip=10mm]{geometry}
\usepackage{xcolor,array,booktabs,longtable,tabularx}
\usepackage{enumitem,fancyhdr}
\usepackage[unicode,hidelinks]{hyperref}
\definecolor{ink}{HTML}{000000}
\definecolor{muted}{HTML}{4C5558}
\setlength{\parindent}{0pt}
\setlength{\parskip}{3.5pt}
\AtBeginDocument{\fontsize{11}{14.3}\selectfont}
\setlength{\tabcolsep}{4pt}
\renewcommand{\arraystretch}{1.10}
\setlist[itemize]{leftmargin=1.3em,label=\textbullet,itemsep=2pt,topsep=3pt,parsep=0pt}
\setlist[enumerate]{leftmargin=1.8em,itemsep=3pt,topsep=4pt,parsep=0pt}
\widowpenalty=10000
\clubpenalty=10000
\emergencystretch=1em
\raggedbottom
\pagestyle{fancy}
\fancyhf{}
\providecommand{\runningtitle}{Old English}
\fancyhead[LE,RO]{\small\itshape \runningtitle}
\fancyhead[LO,RE]{\small Old English}
\fancyfoot[LE,RO]{\small\thepage}
\renewcommand{\headrulewidth}{0.3pt}
\newcommand{\booktitle}[2]{
  {\fontsize{9}{11}\selectfont\color{muted}#1\par}
  \vspace{3pt}
  {\fontsize{20}{24}\selectfont\color{ink}#2\par}
  \vspace{6pt}
}
\newcommand{\source}[1]{\par{\footnotesize\color{muted}#1\par}}
\newcommand{\readingtext}[1]{
  \par\addvspace{5pt}
  \begingroup\leftskip=1.5em
    \fontsize{12}{16}\selectfont #1\par
  \endgroup\addvspace{5pt}
}
\makeatletter
\renewcommand{\section}{\@startsection{section}{1}{\z@}
  {-2.4ex \@plus -0.6ex \@minus -0.2ex}{0.8ex \@plus 0.2ex}
  {\normalfont\fontsize{12}{15}\selectfont\bfseries\color{ink}}}
\renewcommand{\subsection}{\@startsection{subsection}{2}{\z@}
  {-1.8ex \@plus -0.5ex \@minus -0.2ex}{0.5ex \@plus 0.2ex}
  {\normalfont\fontsize{11}{14}\selectfont\bfseries\color{ink}}}
\makeatother

\hyphenpenalty=10000
\exhyphenpenalty=10000
\emergencystretch=3em
\newcommand{\itempara}[2]{\par\addvspace{3pt}\noindent\hangindent=1.8em\hangafter=1\makebox[1.8em][l]{#1.}#2\par}
\newcommand{\reserve}[1]{\par\penalty-100\begingroup\dimen0=\pagegoal\advance\dimen0 by-\pagetotal\ifdim\dimen0<#1\newpage\fi\endgroup}
\begin{document}
\raggedright
\booktitle{UNIT 4}{Strong verbs I–III}

% mdparagraph:P001:begin
Strong verbs form their preterite chiefly by changing the stem vowel. This unit gives the complete selected present and preterite indicative of Classes I–III, with the stem changes needed to use them. The noun, pronoun, adjective and weak-verb systems of Units 1–3 remain in use.
% mdparagraph:P001:end

% mdparagraph:P002:begin
In the tables, \textbf{1}, \textbf{2} and \textbf{3} denote first, second and third person; \textbf{sg.} and \textbf{pl.} denote singular and plural. A column labelled \textbf{1–3 pl.} gives the same form for all three plural persons. Dual pronoun subjects also take these plural forms. Use the selected forms when requested; the documented alternatives remain valid Old English. Macrons and the distinction between e and ę are retained.
% mdparagraph:P002:end

\reserve{7\baselineskip}
% mdparagraph:P003:begin
\section*{1. Principal parts and finite stems}
% mdparagraph:P003:end

% teaching: NEW-U04-GRAM-STRONG-WEAK-CONTRAST
% teaching: NEW-U04-GRAM-PRINCIPAL-PARTS-REFERENCE
% teaching: NEW-U04-GRAM-INDICATIVE-SCOPE
% mdparagraph:P004:begin
The weak verbs already learned add a dental element in the preterite. A strong verb instead has a series of stems: the vowel of its present, singular preterite and plural preterite may differ. The \textbf{four principal parts} identify this series. They are conventionally listed as infinitive, singular preterite, plural preterite and past participle.
% mdparagraph:P004:end

% table: U04-PRINCIPAL-PART-POSITIONS
% mdtable:T01:begin
\par\begingroup
\fontsize{10.5}{13.2}\selectfont
\setlength{\LTleft}{0pt}\setlength{\LTright}{\fill}
\noindent
\begin{tabular}{@{}>{\raggedright\arraybackslash}p{33.37mm}>{\raggedright\arraybackslash}p{87.82mm}@{}}
\toprule
\textbf{Position} & \textbf{Function} \\
\midrule
1. Infinitive & Names the verb and identifies the present stem. \\
2. Singular preterite & Identifies the ordinary first/\allowbreak{}third singular preterite stem. \\
3. Plural preterite & Identifies the stem used in all plural persons and in the second singular. \\
4. Past participle & Completes the cited stem series; its syntax is not yet required. \\
\bottomrule
\end{tabular}
\par\endgroup

% mdtable:T01:end
% mdparagraph:P005:begin
Thus drīfan, drāf, drifon, drifen display the Class I sequence ī : ā : i : i. Drāf and drifon are finite indicative forms. The infinitive and participle columns are \textbf{supplied citations}: they may be consulted, and their independent recall or use in infinitive, participial, perfect or passive constructions is not required here. A principal-part entry identifies a stem; the complete tables below identify the person and number.
% mdparagraph:P005:end

% mdparagraph:P006:begin
Present indicative forms state present actions or conditions; preterite indicative forms state past actions or conditions. Every finite table in this unit covers those two tenses and all three persons in singular and plural. The two rare reference verbs in §5 are for recognition only.
% mdparagraph:P006:end

\reserve{7\baselineskip}
% mdparagraph:P007:begin
\section*{2. Endings and stem changes}
% mdparagraph:P007:end

% teaching: NEW-U04-GRAM-PRESENT-ENDINGS
% teaching: NEW-U04-GRAM-PRETERITE-ENDINGS
% table: U04-ENDINGS
% mdtable:T02:begin
\par\begingroup
\fontsize{10.5}{13.2}\selectfont
\setlength{\LTleft}{0pt}\setlength{\LTright}{\fill}
\noindent
\begin{tabular}{@{}>{\raggedright\arraybackslash}p{26.56mm}>{\raggedright\arraybackslash}p{41.22mm}>{\raggedright\arraybackslash}p{50.59mm}@{}}
\toprule
\textbf{Person and number} & \textbf{Present} & \textbf{Preterite} \\
\midrule
1 singular & -e & singular preterite stem; no added ending \\
2 singular & -st, with the listed stem changes & plural preterite stem + -e \\
3 singular & -ð, with the listed stem changes & singular preterite stem; no added ending \\
1, 2, 3 plural & -að & plural preterite stem + -on \\
\bottomrule
\end{tabular}
\par\endgroup

% mdtable:T02:end
% mdparagraph:P008:begin
Contracted verbs have the special present forms printed in their tables. The strong \textbf{second singular preterite} uses the plural stem, not the ordinary first/\allowbreak{}third singular stem: drife belongs with drifon, while drāf is the first/\allowbreak{}third singular. Findan has an additional third-singular exception, given in §4. The indicative tables keep every person explicit.
% mdparagraph:P008:end

% teaching: NEW-U04-GRAM-PRESENT-CONSONANTS
% teaching: NEW-U04-GRAM-GRAMMATICAL-CHANGE
% mdparagraph:P009:begin
The second and third singular present often lose the vowel of older -est and -eð endings. The resulting consonant groups change. Dental d or t before ð gives tt, sometimes written t; d before st gives tst. Stem-final s combines with either ending as st. A stem-final ð may disappear before st, and doubled ð may be reduced. Doubled consonants may also simplify before a consonantal ending. The selected forms below show the result for each verb.
% mdparagraph:P009:end

% mdparagraph:P010:begin
Some verbs change a consonant between their principal-part stems: snīðan and sēoðan have d in their plural-based preterites, while cēosan has r. Such alternation belongs to the particular verb and stem position. Likewise, g/\allowbreak{}h alternatives occur in the stated environments after a long vowel or a liquid; they are not interchangeable throughout a paradigm.
% mdparagraph:P010:end

% mdparagraph:P011:begin
Full -est/\allowbreak{}-eð forms also occur. Their vowels may retain an older change or follow another part of the paradigm by analogy. The notes list documented alternatives; the selected tables do not license automatic manufacture of every possible full-ending spelling.
% mdparagraph:P011:end

\reserve{7\baselineskip}
% mdparagraph:P012:begin
\section*{3. Classes I and II}
% mdparagraph:P012:end

\reserve{7\baselineskip}
% mdparagraph:P013:begin
\subsection*{Class I}
% mdparagraph:P013:end

% teaching: NEW-U04-GRAM-CLASS-I
% mdparagraph:P014:begin
The ordinary Class I series is \textbf{ī : ā : i : i}. The long ī remains in the selected present stem; the preterite singular has ā, and the plural-based preterite has short i. The fourth vowel belongs to the supplied participle citation.
% mdparagraph:P014:end

% teaching: NEW-U04-VOC-drifan
% teaching: NEW-U04-CIT-drifan
% teaching: NEW-U04-VOC-ridan
% teaching: NEW-U04-CIT-ridan
% teaching: NEW-U04-VOC-writan
% teaching: NEW-U04-CIT-writan
% teaching: NEW-U04-VOC-risan
% teaching: NEW-U04-CIT-risan
% teaching: NEW-U04-VOC-snithan
% teaching: NEW-U04-CIT-snithan
% teaching: NEW-U04-VOC-wreon
% teaching: NEW-U04-CIT-wreon
% table: U04-PARTS-I
% mdtable:T03:begin
\par\begingroup
\fontsize{10.5}{13.2}\selectfont
\setlength{\LTleft}{0pt}\setlength{\LTright}{\fill}
\noindent
\begin{tabular}{@{}>{\raggedright\arraybackslash}p{23.14mm}>{\raggedright\arraybackslash}p{21.04mm}>{\raggedright\arraybackslash}p{22.72mm}>{\raggedright\arraybackslash}p{22.72mm}>{\raggedright\arraybackslash}p{23.14mm}@{}}
\toprule
\textbf{Infinitive¹} & \textbf{Meaning} & \textbf{Preterite sg.} & \textbf{Preterite pl.} & \textbf{Participle¹} \\
\midrule
drīfan & drive & drāf & drifon & drifen \\
rīdan & ride & rād & ridon & riden \\
wrītan & write & wrāt & writon & writen \\
rīsan & rise & rās & rison & risen \\
snīðan & cut & snāð & snidon & sniden \\
wrēon & cover; clothe & wrāh & wrigon & wrigen \\
\bottomrule
\end{tabular}
\par\endgroup

% mdtable:T03:end
% mdparagraph:P015:begin
¹ The infinitive and participle are supplied citations in every principal-part table. The two middle columns identify finite preterite stems; their person assignments are given below.
% mdparagraph:P015:end

% teaching: NEW-U04-PAR-drifan
% teaching: NEW-U04-PAR-ridan
% teaching: NEW-U04-PAR-writan
% teaching: NEW-U04-PAR-risan
% teaching: NEW-U04-PAR-snithan
% teaching: NEW-U04-PAR-wreon
\reserve{11\baselineskip}
% mdparagraph:P016:begin
\textbf{Present indicative.}
% mdparagraph:P016:end

% table: U04-PRESENT-I
% mdtable:T04:begin
\par\begingroup
\fontsize{10.5}{13.2}\selectfont
\setlength{\LTleft}{0pt}\setlength{\LTright}{\fill}
\noindent
\begin{tabular}{@{}>{\raggedright\arraybackslash}p{22.55mm}>{\raggedright\arraybackslash}p{22.55mm}>{\raggedright\arraybackslash}p{22.55mm}>{\raggedright\arraybackslash}p{22.55mm}>{\raggedright\arraybackslash}p{22.55mm}@{}}
\toprule
\textbf{Verb} & \textbf{1 sg.} & \textbf{2 sg.} & \textbf{3 sg.} & \textbf{1–3 pl.} \\
\midrule
drīfan & drīfe & drīfst & drīfð & drīfað \\
rīdan & rīde & rītst & rītt & rīdað \\
wrītan & wrīte & wrītst & wrītt & wrītað \\
rīsan & rīse & rīst & rīst & rīsað \\
snīðan & snīðe & snīst & snīð & snīðað \\
wrēon & wrēo & wrīhst & wrīhð & wrēoð \\
\bottomrule
\end{tabular}
\par\endgroup

% mdtable:T04:end
\reserve{11\baselineskip}
% mdparagraph:P017:begin
\textbf{Preterite indicative.}
% mdparagraph:P017:end

% table: U04-PRETERITE-I
% mdtable:T05:begin
\par\begingroup
\fontsize{10.5}{13.2}\selectfont
\setlength{\LTleft}{0pt}\setlength{\LTright}{\fill}
\noindent
\begin{tabular}{@{}>{\raggedright\arraybackslash}p{22.55mm}>{\raggedright\arraybackslash}p{22.55mm}>{\raggedright\arraybackslash}p{22.55mm}>{\raggedright\arraybackslash}p{22.55mm}>{\raggedright\arraybackslash}p{22.55mm}@{}}
\toprule
\textbf{Verb} & \textbf{1 sg.} & \textbf{2 sg.} & \textbf{3 sg.} & \textbf{1–3 pl.} \\
\midrule
drīfan & drāf & drife & drāf & drifon \\
rīdan & rād & ride & rād & ridon \\
wrītan & wrāt & write & wrāt & writon \\
rīsan & rās & rise & rās & rison \\
snīðan & snāð & snide & snāð & snidon \\
wrēon & wrāh & wrige & wrāh & wrigon \\
\bottomrule
\end{tabular}
\par\endgroup

% mdtable:T05:end
% mdparagraph:P018:begin
In rīdan and wrītan, the selected present second singular ends in -tst and the third in -tt. Third-singular rīt and wrīt are documented alternatives. Rīsan has rīst in both second and third singular; rīsð is another third-singular spelling. Snīðan has selected snīst and snīð, with snīðð also valid in the third singular. Its preterite switches to d in snide and snidon, while the first/\allowbreak{}third singular snāð retains ð.
% mdparagraph:P018:end

% teaching: NEW-U04-GRAM-CONTRACT-I
% mdparagraph:P019:begin
Wrēon has a contracted present: wrēo, wrīhst, wrīhð, wrēoð. Its preterite uses wrāh beside plural-based wrige and wrigon. The h/\allowbreak{}g contrast is part of this pattern. The table selects its Class I series; documented analogical Class II forms are not mixed into that selection. See source note 2.
% mdparagraph:P019:end

\reserve{7\baselineskip}
% mdparagraph:P020:begin
\subsection*{Class II}
% mdparagraph:P020:end

% teaching: NEW-U04-GRAM-CLASS-II
% teaching: NEW-U04-GRAM-PRESENT-MUTATION
% mdparagraph:P021:begin
The ordinary Class II series is \textbf{ēo or ū : ēa : u : o}. In the selected second and third singular present, ēo changes to īe and ū to ȳ. The first singular and plural retain the present vowel shown in the infinitive citation. Cook also records ī beside īe; the tables consistently select īe.
% mdparagraph:P021:end

% teaching: NEW-U04-VOC-beodan
% teaching: NEW-U04-CIT-beodan
% teaching: NEW-U04-VOC-bugan
% teaching: NEW-U04-CIT-bugan
% teaching: NEW-U04-VOC-ceosan
% teaching: NEW-U04-CIT-ceosan
% teaching: NEW-U04-VOC-seothan
% teaching: NEW-U04-CIT-seothan
% teaching: NEW-U04-VOC-dreogan
% teaching: NEW-U04-CIT-dreogan
% teaching: NEW-U04-VOC-teon
% teaching: NEW-U04-CIT-teon
% table: U04-PARTS-II
% mdtable:T06:begin
\par\begingroup
\fontsize{10.5}{13.2}\selectfont
\setlength{\LTleft}{0pt}\setlength{\LTright}{\fill}
\noindent
\begin{tabular}{@{}>{\raggedright\arraybackslash}p{23.14mm}>{\raggedright\arraybackslash}p{21.04mm}>{\raggedright\arraybackslash}p{22.72mm}>{\raggedright\arraybackslash}p{22.72mm}>{\raggedright\arraybackslash}p{23.14mm}@{}}
\toprule
\textbf{Infinitive¹} & \textbf{Meaning} & \textbf{Preterite sg.} & \textbf{Preterite pl.} & \textbf{Participle¹} \\
\midrule
bēodan & offer & bēad & budon & boden \\
būgan & bend & bēag & bugon & bogen \\
cēosan & choose & cēas & curon & coren \\
sēoðan & boil & sēað & sudon & soden \\
drēogan & endure & drēag & drugon & drogen \\
tēon & draw; pull & tēah & tugon & togen \\
\bottomrule
\end{tabular}
\par\endgroup

% mdtable:T06:end
% teaching: NEW-U04-PAR-beodan
% teaching: NEW-U04-PAR-bugan
% teaching: NEW-U04-PAR-ceosan
% teaching: NEW-U04-PAR-seothan
% teaching: NEW-U04-PAR-dreogan
% teaching: NEW-U04-PAR-teon
\reserve{11\baselineskip}
% mdparagraph:P022:begin
\textbf{Present indicative.}
% mdparagraph:P022:end

% table: U04-PRESENT-II
% mdtable:T07:begin
\par\begingroup
\fontsize{10.5}{13.2}\selectfont
\setlength{\LTleft}{0pt}\setlength{\LTright}{\fill}
\noindent
\begin{tabular}{@{}>{\raggedright\arraybackslash}p{22.55mm}>{\raggedright\arraybackslash}p{22.55mm}>{\raggedright\arraybackslash}p{22.55mm}>{\raggedright\arraybackslash}p{22.55mm}>{\raggedright\arraybackslash}p{22.55mm}@{}}
\toprule
\textbf{Verb} & \textbf{1 sg.} & \textbf{2 sg.} & \textbf{3 sg.} & \textbf{1–3 pl.} \\
\midrule
bēodan & bēode & bīetst & bīett & bēodað \\
būgan & būge & bȳgst & bȳgð & būgað \\
cēosan & cēose & cīest & cīest & cēosað \\
sēoðan & sēoðe & sīest & sīeð & sēoðað \\
drēogan & drēoge & drīegst & drīegð & drēogað \\
tēon & tēo & tīehst & tīehð & tēoð \\
\bottomrule
\end{tabular}
\par\endgroup

% mdtable:T07:end
\reserve{11\baselineskip}
% mdparagraph:P023:begin
\textbf{Preterite indicative.}
% mdparagraph:P023:end

% table: U04-PRETERITE-II
% mdtable:T08:begin
\par\begingroup
\fontsize{10.5}{13.2}\selectfont
\setlength{\LTleft}{0pt}\setlength{\LTright}{\fill}
\noindent
\begin{tabular}{@{}>{\raggedright\arraybackslash}p{22.55mm}>{\raggedright\arraybackslash}p{22.55mm}>{\raggedright\arraybackslash}p{22.55mm}>{\raggedright\arraybackslash}p{22.55mm}>{\raggedright\arraybackslash}p{22.55mm}@{}}
\toprule
\textbf{Verb} & \textbf{1 sg.} & \textbf{2 sg.} & \textbf{3 sg.} & \textbf{1–3 pl.} \\
\midrule
bēodan & bēad & bude & bēad & budon \\
būgan & bēag & buge & bēag & bugon \\
cēosan & cēas & cure & cēas & curon \\
sēoðan & sēað & sude & sēað & sudon \\
drēogan & drēag & druge & drēag & drugon \\
tēon & tēah & tuge & tēah & tugon \\
\bottomrule
\end{tabular}
\par\endgroup

% mdtable:T08:end
% mdparagraph:P024:begin
Bēodan combines the vowel change with dental endings: bīetst and bīett; third-singular bīet is also valid. Būgan has selected bȳgst/\allowbreak{}bȳgð and the documented alternatives bȳhst/\allowbreak{}bȳhð. Its singular preterite bēag may also be written bēah.
% mdparagraph:P024:end

% mdparagraph:P025:begin
Cēosan has cīest in both second and third singular present, cēas in the first/\allowbreak{}third singular preterite, and r in cure and curon. Sēoðan has sīest and sīeð in the present, sēað in the first/\allowbreak{}third singular preterite, and d in sude and sudon. The fourth principal parts retain the corresponding coren and soden stems.
% mdparagraph:P025:end

% mdparagraph:P026:begin
Drēogan selects drīegst/\allowbreak{}drīegð in the present; drīehst/\allowbreak{}drīehð are documented alternatives. Its first/\allowbreak{}third singular preterite may be drēah beside selected drēag. Keep those g/\allowbreak{}h choices distinct from the s/\allowbreak{}r and ð/\allowbreak{}d alternations of the preceding verbs.
% mdparagraph:P026:end

% teaching: NEW-U04-GRAM-CONTRACT-II
% mdparagraph:P027:begin
Tēon “draw, pull” is the contracted Class II verb here. Its present is tēo, tīehst, tīehð, tēoð; the preterite has tēah beside tuge and tugon. Its similar infinitive ending does not give it the Class I preterite of wrēon. The meaning identifies the assigned verb; no other homographic verb is included.
% mdparagraph:P027:end

\reserve{7\baselineskip}
% mdparagraph:P028:begin
\section*{4. Class III}
% mdparagraph:P028:end

% teaching: NEW-U04-GRAM-CLASS-III-NASAL
% teaching: NEW-U04-GRAM-CLASS-III-EL-IEL
% teaching: NEW-U04-GRAM-CLASS-III-EOR
% teaching: NEW-U04-GRAM-CLASS-III-METATHESIS
% mdparagraph:P029:begin
Class III contains several related stem patterns. Its vowels in the following series are \textbf{short}. The present second/\allowbreak{}third singular may differ from the infinitive vowel, and the consonant changes described in §2 still apply.
% mdparagraph:P029:end

% table: U04-III-SERIES
% mdtable:T09:begin
\par\begingroup
\fontsize{10.5}{13.2}\selectfont
\setlength{\LTleft}{0pt}\setlength{\LTright}{\fill}
\noindent
\begin{tabular}{@{}>{\raggedright\arraybackslash}p{30.49mm}>{\raggedright\arraybackslash}p{34.43mm}>{\raggedright\arraybackslash}p{53.45mm}@{}}
\toprule
\textbf{Family} & \textbf{Four-part vowel series} & \textbf{Assigned verbs} \\
\midrule
Nasal + consonant & i : a : u : u & bindan, winnan, findan, drincan, singan \\
el & e : ea : u : o & helpan \\
iel & ie : ea : u : o & gieldan \\
eor & eo : ea : u : o & weorpan, weorðan, forweorðan \\
Metathesis & e : æ : u : o & berstan, tōberstan \\
\bottomrule
\end{tabular}
\par\endgroup

% mdtable:T09:end
% mdparagraph:P030:begin
Metathesis is a change in the order of sounds. In the berstan family, r and the neighbouring vowel have changed order. Its selected series is berstan, bærst, burston, borsten; it is not obtained by substituting the weorpan vowel in every cell.
% mdparagraph:P030:end

% teaching: NEW-U04-VOC-bindan
% teaching: NEW-U04-CIT-bindan
% teaching: NEW-U04-VOC-helpan
% teaching: NEW-U04-CIT-helpan
% teaching: NEW-U04-VOC-gieldan
% teaching: NEW-U04-CIT-gieldan
% teaching: NEW-U04-VOC-weorpan
% teaching: NEW-U04-CIT-weorpan
% teaching: NEW-U04-VOC-berstan
% teaching: NEW-U04-CIT-berstan
% teaching: NEW-U04-VOC-weorthan
% teaching: NEW-U04-CIT-weorthan
% teaching: NEW-U04-VOC-winnan
% teaching: NEW-U04-CIT-winnan
% teaching: NEW-U04-VOC-findan
% teaching: NEW-U04-CIT-findan
% teaching: NEW-U04-VOC-drincan
% teaching: NEW-U04-CIT-drincan
% teaching: NEW-U04-VOC-singan
% teaching: NEW-U04-CIT-singan
% teaching: NEW-U04-VOC-toberstan
% teaching: NEW-U04-CIT-toberstan
% teaching: NEW-U04-VOC-forweorthan
% teaching: NEW-U04-CIT-forweorthan
% table: U04-PARTS-III
% mdtable:T10:begin
\par\begingroup
\fontsize{10.5}{13.2}\selectfont
\setlength{\LTleft}{0pt}\setlength{\LTright}{\fill}
\noindent
\begin{tabular}{@{}>{\raggedright\arraybackslash}p{21.22mm}>{\raggedright\arraybackslash}p{28.65mm}>{\raggedright\arraybackslash}p{20.83mm}>{\raggedright\arraybackslash}p{20.83mm}>{\raggedright\arraybackslash}p{21.22mm}@{}}
\toprule
\textbf{Infinitive¹} & \textbf{Meaning} & \textbf{Preterite sg.} & \textbf{Preterite pl.} & \textbf{Participle¹} \\
\midrule
bindan & bind; tie & band & bundon & bunden \\
helpan & help & healp & hulpon & holpen \\
gieldan & yield & geald & guldon & golden \\
weorpan & throw & wearp & wurpon & worpen \\
berstan & burst & bærst & burston & borsten \\
weorðan & become & wearð & wurdon & worden \\
winnan & struggle; toil & wann & wunnon & wunnen \\
findan & find & fand & fundon & funden \\
drincan & drink & dranc & druncon & druncen \\
singan & sing & sang & sungon & sungen \\
tōberstan & break up; go to pieces & tōbærst & tōburston & tōborsten \\
forweorðan & perish & forwearð & forwurdon & forworden \\
\bottomrule
\end{tabular}
\par\endgroup

% mdtable:T10:end
% teaching: NEW-U04-PAR-bindan
% teaching: NEW-U04-PAR-helpan
% teaching: NEW-U04-PAR-gieldan
% teaching: NEW-U04-PAR-weorpan
% teaching: NEW-U04-PAR-berstan
% teaching: NEW-U04-PAR-weorthan
% teaching: NEW-U04-PAR-winnan
% teaching: NEW-U04-PAR-findan
% teaching: NEW-U04-PAR-drincan
% teaching: NEW-U04-PAR-singan
% teaching: NEW-U04-PAR-toberstan
% teaching: NEW-U04-PAR-forweorthan
\reserve{17\baselineskip}
% mdparagraph:P031:begin
\textbf{Present indicative.}
% mdparagraph:P031:end

% table: U04-PRESENT-III
% mdtable:T11:begin
\par\begingroup
\fontsize{10.5}{13.2}\selectfont
\setlength{\LTleft}{0pt}\setlength{\LTright}{\fill}
\noindent
\begin{tabular}{@{}>{\raggedright\arraybackslash}p{22.55mm}>{\raggedright\arraybackslash}p{22.55mm}>{\raggedright\arraybackslash}p{22.55mm}>{\raggedright\arraybackslash}p{22.55mm}>{\raggedright\arraybackslash}p{22.55mm}@{}}
\toprule
\textbf{Verb} & \textbf{1 sg.} & \textbf{2 sg.} & \textbf{3 sg.} & \textbf{1–3 pl.} \\
\midrule
bindan & binde & bintst & bint & bindað \\
helpan & helpe & hilpst & hilpð & helpað \\
gieldan & gielde & gieltst & gielt & gieldað \\
weorpan & weorpe & wierpst & wierpð & weorpað \\
berstan & berste & bierst & bierst & berstað \\
weorðan & weorðe & wierst & wierð & weorðað \\
winnan & winne & winst & winð & winnað \\
findan & finde & fintst & fint & findað \\
drincan & drince & drincst & drincð & drincað \\
singan & singe & singst & singð & singað \\
tōberstan & tōberste & tōbierst & tōbierst & tōberstað \\
forweorðan & forweorðe & forwierst & forwierð & forweorðað \\
\bottomrule
\end{tabular}
\par\endgroup

% mdtable:T11:end
\reserve{17\baselineskip}
% mdparagraph:P032:begin
\textbf{Preterite indicative.}
% mdparagraph:P032:end

% table: U04-PRETERITE-III
% mdtable:T12:begin
\par\begingroup
\fontsize{10.5}{13.2}\selectfont
\setlength{\LTleft}{0pt}\setlength{\LTright}{\fill}
\noindent
\begin{tabular}{@{}>{\raggedright\arraybackslash}p{22.55mm}>{\raggedright\arraybackslash}p{22.55mm}>{\raggedright\arraybackslash}p{22.55mm}>{\raggedright\arraybackslash}p{22.55mm}>{\raggedright\arraybackslash}p{22.55mm}@{}}
\toprule
\textbf{Verb} & \textbf{1 sg.} & \textbf{2 sg.} & \textbf{3 sg.} & \textbf{1–3 pl.} \\
\midrule
bindan & band & bunde & band & bundon \\
helpan & healp & hulpe & healp & hulpon \\
gieldan & geald & gulde & geald & guldon \\
weorpan & wearp & wurpe & wearp & wurpon \\
berstan & bærst & burste & bærst & burston \\
weorðan & wearð & wurde & wearð & wurdon \\
winnan & wann & wunne & wann & wunnon \\
findan & fand & funde & funde & fundon \\
drincan & dranc & drunce & dranc & druncon \\
singan & sang & sunge & sang & sungon \\
tōberstan & tōbærst & tōburste & tōbærst & tōburston \\
forweorðan & forwearð & forwurde & forwearð & forwurdon \\
\bottomrule
\end{tabular}
\par\endgroup

% mdtable:T12:end
% mdparagraph:P033:begin
Bindan has bintst/\allowbreak{}bint in the selected present. Bindest or bindst may occur beside bintst, and bindeð or bindð beside bint. Winnan reduces nn before the singular consonantal endings, giving winst and winð, but retains nn in winne, winnað and the preterite stems. Drincan and singan follow the nasal vowel series without those dental or doubled-n changes.
% mdparagraph:P033:end

% mdparagraph:P034:begin
Helpan has hilpst/\allowbreak{}hilpð; gieldan has gieltst/\allowbreak{}gielt; weorpan has wierpst/\allowbreak{}wierpð. The vowel and the final consonant group both matter. Berstan has bierst in both second and third singular present. Weorðan has selected wierst, with wierðst also valid, and third-singular wierð. Its first/\allowbreak{}third singular preterite wearð retains ð, while wurde and wurdon have d.
% mdparagraph:P034:end

% teaching: NEW-U04-GRAM-FINDAN-EXCEPTION
% mdparagraph:P035:begin
Findan has selected present fintst/\allowbreak{}fint. Its first singular preterite is \textbf{fand}, its second singular \textbf{funde}, and its usual third singular \textbf{funde}. Thus the two singular funde forms are distinguished by person or context, not by spelling. Third-singular fand is valid; first-singular funde is also documented. A request for the selected pattern requires the forms in the table, but these alternatives are not grammatical errors.
% mdparagraph:P035:end

% teaching: NEW-U04-GRAM-PREFIXED-LEXEMES
% teaching: NEW-U04-GRAM-PREFIXED-INTRANSITIVE
% mdparagraph:P036:begin
Tōberstan “break up, go to pieces” follows berstan; forweorðan “perish” follows weorðan. Both are used intransitively in the source reading: the ship breaks up and the companions perish. Their complete forms and whole-word meanings are learned separately. The prefixes do not by themselves supply a general rule for creating other verbs or predicting their meanings. Forwierðst is valid beside selected forwierst.
% mdparagraph:P036:end

\reserve{7\baselineskip}
% mdparagraph:P037:begin
\section*{5. Complete recognition references}
% mdparagraph:P037:end

% teaching: NEW-U04-GRAM-RARE-REFERENCE
% mdparagraph:P038:begin
The following two exceptional systems are supplied as complete \textbf{recognition references}. Consult the tables for the reference exercise; independent production and unintroduced object constructions are not required.
% mdparagraph:P038:end

% teaching: NEW-U04-VOC-bregdan
% teaching: NEW-U04-CIT-bregdan
% teaching: NEW-U04-VOC-feolan
% teaching: NEW-U04-CIT-feolan
% table: U04-PARTS-REFERENCE
% mdtable:T13:begin
\par\begingroup
\fontsize{10.5}{13.2}\selectfont
\setlength{\LTleft}{0pt}\setlength{\LTright}{\fill}
\noindent
\begin{tabular}{@{}>{\raggedright\arraybackslash}p{22.61mm}>{\raggedright\arraybackslash}p{23.13mm}>{\raggedright\arraybackslash}p{22.2mm}>{\raggedright\arraybackslash}p{22.2mm}>{\raggedright\arraybackslash}p{22.61mm}@{}}
\toprule
\textbf{Infinitive¹} & \textbf{Meaning} & \textbf{Preterite sg.} & \textbf{Preterite pl.} & \textbf{Participle¹} \\
\midrule
bregdan & brandish; draw & brægd & brugdon & brogden \\
fēolan & penetrate & fealh & fulgon & folen \\
\bottomrule
\end{tabular}
\par\endgroup

% mdtable:T13:end
% teaching: NEW-U04-PAR-bregdan
% teaching: NEW-U04-PAR-feolan
\reserve{7\baselineskip}
% mdparagraph:P039:begin
\textbf{Present indicative.}
% mdparagraph:P039:end

% table: U04-PRESENT-REFERENCE
% mdtable:T14:begin
\par\begingroup
\fontsize{10.5}{13.2}\selectfont
\setlength{\LTleft}{0pt}\setlength{\LTright}{\fill}
\noindent
\begin{tabular}{@{}>{\raggedright\arraybackslash}p{22.55mm}>{\raggedright\arraybackslash}p{22.55mm}>{\raggedright\arraybackslash}p{22.55mm}>{\raggedright\arraybackslash}p{22.55mm}>{\raggedright\arraybackslash}p{22.55mm}@{}}
\toprule
\textbf{Verb} & \textbf{1 sg.} & \textbf{2 sg.} & \textbf{3 sg.} & \textbf{1–3 pl.} \\
\midrule
bregdan & bregde & brītst & brītt & bregdað \\
fēolan & fēole & filhst & filhð & fēolað \\
\bottomrule
\end{tabular}
\par\endgroup

% mdtable:T14:end
\reserve{7\baselineskip}
% mdparagraph:P040:begin
\textbf{Preterite indicative.}
% mdparagraph:P040:end

% table: U04-PRETERITE-REFERENCE
% mdtable:T15:begin
\par\begingroup
\fontsize{10.5}{13.2}\selectfont
\setlength{\LTleft}{0pt}\setlength{\LTright}{\fill}
\noindent
\begin{tabular}{@{}>{\raggedright\arraybackslash}p{22.55mm}>{\raggedright\arraybackslash}p{22.55mm}>{\raggedright\arraybackslash}p{22.55mm}>{\raggedright\arraybackslash}p{22.55mm}>{\raggedright\arraybackslash}p{22.55mm}@{}}
\toprule
\textbf{Verb} & \textbf{1 sg.} & \textbf{2 sg.} & \textbf{3 sg.} & \textbf{1–3 pl.} \\
\midrule
bregdan & brægd & brugde & brægd & brugdon \\
fēolan & fealh & fulge & fealh & fulgon \\
\bottomrule
\end{tabular}
\par\endgroup

% mdtable:T15:end
% mdparagraph:P041:begin
Bregdan loses g and lengthens the vowel in selected brītst/\allowbreak{}brītt; brīt is also valid in the third singular. In its preterite, brǣd may occur beside brægd, and brūdon beside brugdon. The participle citation brogden also has the variant brōden. These are reference alternatives, not additional recall targets.
% mdparagraph:P041:end

% mdparagraph:P042:begin
Fēolan has the exceptional selected present filhst/\allowbreak{}filhð beside fēole/\allowbreak{}fēolað. Its preterite uses fealh, fulge and fulgon; fǣlon is a documented plural alternative. The participle citation is Cook’s folen. Wright records folgen and fōlen separately. Preserve these distinctions rather than adding or removing a length mark. This complete reference combines the sources described in note 3.
% mdparagraph:P042:end

\reserve{7\baselineskip}
% mdparagraph:P043:begin
\section*{6. Four nouns}
% mdparagraph:P043:end

% teaching: NEW-U04-GRAM-NEW-NOUNS
% teaching: NEW-U04-VOC-scip
% teaching: NEW-U04-VOC-wind
% teaching: NEW-U04-VOC-strand
% teaching: NEW-U04-VOC-gefera
% table: U04-NOUN-VOCABULARY
% mdtable:T16:begin
\par\begingroup
\fontsize{10.5}{13.2}\selectfont
\setlength{\LTleft}{0pt}\setlength{\LTright}{\fill}
\noindent
\begin{tabular}{@{}>{\raggedright\arraybackslash}p{25.77mm}>{\raggedright\arraybackslash}p{25.77mm}>{\raggedright\arraybackslash}p{25.77mm}>{\raggedright\arraybackslash}p{38.26mm}@{}}
\toprule
\textbf{Noun} & \textbf{Gender} & \textbf{Meaning} & \textbf{Pattern} \\
\midrule
scip & neuter & ship & short neuter, like hof \\
wind & masculine & wind & strong masculine \\
strand & masculine & shore & strong masculine \\
gefēra & masculine & companion & weak masculine \\
\bottomrule
\end{tabular}
\par\endgroup

% mdtable:T16:end
% teaching: NEW-U04-PAR-scip
% teaching: NEW-U04-PAR-wind
% table: U04-NOUN-SCIP-WIND
% mdtable:T17:begin
\par\begingroup
\fontsize{10.5}{13.2}\selectfont
\setlength{\LTleft}{0pt}\setlength{\LTright}{\fill}
\noindent
\begin{tabular}{@{}>{\raggedright\arraybackslash}p{22.55mm}>{\raggedright\arraybackslash}p{22.55mm}>{\raggedright\arraybackslash}p{22.55mm}>{\raggedright\arraybackslash}p{22.55mm}>{\raggedright\arraybackslash}p{22.55mm}@{}}
\toprule
\textbf{Case} & \textbf{scip sg.} & \textbf{scip pl.} & \textbf{wind sg.} & \textbf{wind pl.} \\
\midrule
Nom. & scip & scipu & wind & windas \\
Gen. & scipes & scipa & windes & winda \\
Dat. & scipe & scipum & winde & windum \\
Acc. & scip & scipu & wind & windas \\
\bottomrule
\end{tabular}
\par\endgroup

% mdtable:T17:end
% teaching: NEW-U04-PAR-strand
% teaching: NEW-U04-PAR-gefera
% table: U04-NOUN-STRAND-GEFERA
% mdtable:T18:begin
\par\begingroup
\fontsize{10.5}{13.2}\selectfont
\setlength{\LTleft}{0pt}\setlength{\LTright}{\fill}
\noindent
\begin{tabular}{@{}>{\raggedright\arraybackslash}p{22.55mm}>{\raggedright\arraybackslash}p{22.55mm}>{\raggedright\arraybackslash}p{22.55mm}>{\raggedright\arraybackslash}p{22.55mm}>{\raggedright\arraybackslash}p{22.55mm}@{}}
\toprule
\textbf{Case} & \textbf{strand sg.} & \textbf{strand pl.} & \textbf{gefēra sg.} & \textbf{gefēra pl.} \\
\midrule
Nom. & strand & strandas & gefēra & gefēran \\
Gen. & strandes & stranda & gefēran & gefērena \\
Dat. & strande & strandum & gefēran & gefērum \\
Acc. & strand & strandas & gefēran & gefēran \\
\bottomrule
\end{tabular}
\par\endgroup

% mdtable:T18:end
% mdparagraph:P044:begin
Scip has a short stem syllable and therefore nominative/\allowbreak{}accusative plural scipu. The initial sc does not make the vowel syllable heavy: both consonants precede the vowel. Wind and strand retain their ordinary strong masculine stems. Gefēra has the weak oblique singular and nominative/\allowbreak{}accusative plural gefēran, but genitive plural gefērena and dative plural gefērum. The same written form gefēran can have several analyses; its clause supplies the choice.
% mdparagraph:P044:end

\reserve{7\baselineskip}
% mdparagraph:P045:begin
\section*{7. Ordinary clauses}
% mdparagraph:P045:end

% teaching: NEW-U04-GRAM-DIRECT-OBJECTS
% mdparagraph:P046:begin
The following selected uses take an \textbf{accusative object}: drīfan with the being driven, bindan with the being bound, cēosan with the person or thing chosen, findan with the person or thing found, and drincan with the liquid drunk. The verb’s particular use determines its construction. These object patterns do not follow merely from membership in a strong class.
% mdparagraph:P046:end

% mdparagraph:P047:begin
The eight examples in this section are newly composed teaching material. The other verbs in the tables may be practised as labelled finite forms without assuming additional object or complement constructions.
% mdparagraph:P047:end

% example: NEW-U04-EXAMPLE-01
% mdparagraph:P048:begin
Ic drīfe ðā nīetenu.
% mdparagraph:P048:end

% mdparagraph:P049:begin
I drive the cattle.
% mdparagraph:P049:end

% mdparagraph:P050:begin
Nīetenu is the accusative plural object. The singular first-person subject takes drīfe.
% mdparagraph:P050:end

% example: NEW-U04-EXAMPLE-03
% mdparagraph:P051:begin
Ðū cīest ðæt scip.
% mdparagraph:P051:end

% mdparagraph:P052:begin
You choose the ship.
% mdparagraph:P052:end

% mdparagraph:P053:begin
Scip is the neuter accusative singular object. Cīest is the selected second singular present of cēosan.
% mdparagraph:P053:end

% example: NEW-U04-EXAMPLE-04
% mdparagraph:P054:begin
Se gefēra dranc wæter.
% mdparagraph:P054:end

% mdparagraph:P055:begin
The companion drank water.
% mdparagraph:P055:end

% mdparagraph:P056:begin
The companion is the subject; wæter is the accusative object. Dranc is the third singular preterite.
% mdparagraph:P056:end

% teaching: NEW-U04-GRAM-INTRANSITIVE-ACTIONS
% mdparagraph:P057:begin
Rīdan “ride” and singan “sing” may stand without a nominal object. A familiar adverb may qualify the action. Riding here has its ordinary horseback sense; these clauses introduce no direction or means phrase.
% mdparagraph:P057:end

% example: NEW-U04-EXAMPLE-02
% mdparagraph:P058:begin
Ðū rītst and hē rītt.
% mdparagraph:P058:end

% mdparagraph:P059:begin
You ride, and he rides.
% mdparagraph:P059:end

% mdparagraph:P060:begin
Each singular subject has its own finite verb. Both forms preserve the dental changes of rīdan.
% mdparagraph:P060:end

% example: NEW-U04-EXAMPLE-08
% mdparagraph:P061:begin
Wit sungon wel.
% mdparagraph:P061:end

% mdparagraph:P062:begin
We two sang well.
% mdparagraph:P062:end

% mdparagraph:P063:begin
The dual subject takes the plural preterite sungon. Wel qualifies the singing.
% mdparagraph:P063:end

% teaching: NEW-U04-GRAM-HELPAN-DATIVE
% mdparagraph:P064:begin
The selected construction of \textbf{helpan} puts the recipient of help in the \textbf{dative}. A person or another being helped can occupy this position. The genitive also occurs with this verb; it is valid Old English but differs from the dative construction requested here. The use with birds in the cumulative practice is an ordinary application of the same helping sense.
% mdparagraph:P064:end

% example: NEW-U04-EXAMPLE-05
% mdparagraph:P065:begin
Hē hilpð ðǣm męnn.
% mdparagraph:P065:end

% mdparagraph:P066:begin
He helps the man.
% mdparagraph:P066:end

% mdparagraph:P067:begin
Męnn is dative singular, marked also by ðǣm. It is the recipient of help, not an accusative object.
% mdparagraph:P067:end

% teaching: NEW-U04-GRAM-BECOME-PREDICATE
% mdparagraph:P068:begin
Ordinary \textbf{weorðan} means “become”. Its predicate adjective is nominative and agrees with the subject in gender and number, as a predicate with the copula does. This ordinary use is distinct from the passive expressions previously supplied in readings.
% mdparagraph:P068:end

% example: NEW-U04-EXAMPLE-06
% mdparagraph:P069:begin
Se gefēra wearð glæd.
% mdparagraph:P069:end

% mdparagraph:P070:begin
The companion became glad.
% mdparagraph:P070:end

% mdparagraph:P071:begin
Glæd is strong masculine nominative singular, agreeing with gefēra. Wearð is the third singular preterite of weorðan.
% mdparagraph:P071:end

% teaching: NEW-U04-GRAM-SIMPLE-NEGATION
% mdparagraph:P072:begin
Place \textbf{ne} before the noncontracting strong finite forms used here, just as before the ordinary weak forms already learned. The form of the verb otherwise keeps its person, number and tense. Contracted negative copulas and other negative systems are not part of this rule.
% mdparagraph:P072:end

% example: NEW-U04-EXAMPLE-07
% mdparagraph:P073:begin
Ic ne finde ðæt scip.
% mdparagraph:P073:end

% mdparagraph:P074:begin
I do not find the ship.
% mdparagraph:P074:end

% mdparagraph:P075:begin
Ne negates the first singular present finde. The object remains accusative.
% mdparagraph:P075:end

\reserve{7\baselineskip}
% mdparagraph:P076:begin
\section*{8. Practice}
% mdparagraph:P076:end

% mdparagraph:P077:begin
All exercises in this unit are new teaching material. Principal-part citations and the two reference paradigms remain available where the prompt calls for them. Other production tasks use the selected forms and the ordinary constructions taught above.
% mdparagraph:P077:end

\reserve{7\baselineskip}
% mdparagraph:P078:begin
\subsection*{Principal parts — P1–P4}
% mdparagraph:P078:end

\reserve{5\baselineskip}
\par\begin{minipage}{\linewidth}\raggedright
% mdparagraph:P079:begin
\textbf{P1.}
% mdparagraph:P079:end

% item: NEW-U04-P01
% mdparagraph:P080:begin
The printed entry gives drīfan, drāf, drifon, drifen. Label the four principal-part positions. Which two entries are finite indicative forms, and which two are nonfinite citations?
% mdparagraph:P080:end

\end{minipage}\par\addvspace{4pt}

\par\begin{minipage}{\linewidth}\raggedright
% mdparagraph:P081:begin
\textbf{P2.}
% mdparagraph:P081:end

% item: NEW-U04-P02
% mdparagraph:P082:begin
Using the printed principal-part table, identify the strong class of rīdan, bēodan and bindan. For each verb give the second-person singular preterite and state which principal-part stem it uses.
% mdparagraph:P082:end

\end{minipage}\par\addvspace{4pt}

\par\begin{minipage}{\linewidth}\raggedright
% mdparagraph:P083:begin
\textbf{P3.}
% mdparagraph:P083:end

% item: NEW-U04-P03
% mdparagraph:P084:begin
Give every person and number represented by the finite forms drāf and drifon in the selected indicative paradigm. Explain why the supplied citation drifen is not another person of that paradigm.
% mdparagraph:P084:end

\end{minipage}\par\addvspace{4pt}

\par\begin{minipage}{\linewidth}\raggedright
% mdparagraph:P085:begin
\textbf{P4.}
% mdparagraph:P085:end

% item: NEW-U04-P04
% mdparagraph:P086:begin
Compare ðū ride and ðū hīerdest. Give the person, number and tense of each verb, and explain the different formation of their preterites.
% mdparagraph:P086:end

\end{minipage}\par\addvspace{4pt}

\reserve{7\baselineskip}
% mdparagraph:P087:begin
\subsection*{Class I — A1–A6}
% mdparagraph:P087:end

\reserve{5\baselineskip}
\par\begin{minipage}{\linewidth}\raggedright
% mdparagraph:P088:begin
\textbf{A1.}
% mdparagraph:P088:end

% item: NEW-U04-A01
% mdparagraph:P089:begin
Give the complete selected present and preterite indicative of drīfan. Label all three persons in the singular and plural even where forms coincide.
% mdparagraph:P089:end

\end{minipage}\par\addvspace{4pt}

\par\begin{minipage}{\linewidth}\raggedright
% mdparagraph:P090:begin
\textbf{A2.}
% mdparagraph:P090:end

% item: NEW-U04-A02
% mdparagraph:P091:begin
For rīdan and wrītan, give the first-, second- and third-person singular present, then the first- and second-person singular preterite and the plural preterite. Explain the selected present dental forms.
% mdparagraph:P091:end

\end{minipage}\par\addvspace{4pt}

\par\begin{minipage}{\linewidth}\raggedright
% mdparagraph:P092:begin
\textbf{A3.}
% mdparagraph:P092:end

% item: NEW-U04-A03
% mdparagraph:P093:begin
For rīsan and snīðan, give the second- and third-person singular present and the second-person singular and plural preterite. Identify the consonant change in snīðan.
% mdparagraph:P093:end

\end{minipage}\par\addvspace{4pt}

\par\begin{minipage}{\linewidth}\raggedright
% mdparagraph:P094:begin
\textbf{A4.}
% mdparagraph:P094:end

% item: NEW-U04-A04
% mdparagraph:P095:begin
Give the complete selected present and preterite indicative of wrēon. Keep its contracted present forms distinct from its preterite stems.
% mdparagraph:P095:end

\end{minipage}\par\addvspace{4pt}

\par\begin{minipage}{\linewidth}\raggedright
% mdparagraph:P096:begin
\textbf{A5.}
% mdparagraph:P096:end

% item: NEW-U04-A05
% mdparagraph:P097:begin
Translate “I rode. You rode. He rode. We rode.” Address one person in the second sentence. Then give the four corresponding present-tense sentences.
% mdparagraph:P097:end

\end{minipage}\par\addvspace{4pt}

\par\begin{minipage}{\linewidth}\raggedright
% mdparagraph:P098:begin
\textbf{A6.}
% mdparagraph:P098:end

% item: NEW-U04-A06
% mdparagraph:P099:begin
With drīfan, give the present forms agreeing with ic, wē, ðū and hē. Then give the corresponding preterite forms. Keep the person and number of each subject unchanged.
% mdparagraph:P099:end

\end{minipage}\par\addvspace{4pt}

\reserve{7\baselineskip}
% mdparagraph:P100:begin
\subsection*{Class II — B1–B6}
% mdparagraph:P100:end

\reserve{5\baselineskip}
\par\begin{minipage}{\linewidth}\raggedright
% mdparagraph:P101:begin
\textbf{B1.}
% mdparagraph:P101:end

% item: NEW-U04-B01
% mdparagraph:P102:begin
Give the complete selected present and preterite indicative of bēodan. Identify the present forms whose vowel differs from that of the infinitive citation.
% mdparagraph:P102:end

\end{minipage}\par\addvspace{4pt}

\par\begin{minipage}{\linewidth}\raggedright
% mdparagraph:P103:begin
\textbf{B2.}
% mdparagraph:P103:end

% item: NEW-U04-B02
% mdparagraph:P104:begin
For būgan, give the first-, second- and third-person singular present and the second-person singular and plural preterite. Contrast its present vowel change with bēodan.
% mdparagraph:P104:end

\end{minipage}\par\addvspace{4pt}

\par\begin{minipage}{\linewidth}\raggedright
% mdparagraph:P105:begin
\textbf{B3.}
% mdparagraph:P105:end

% item: NEW-U04-B03
% mdparagraph:P106:begin
For cēosan, give the second- and third-person singular present, the first- and second-person singular preterite and the plural preterite. Explain where the selected s and r occur.
% mdparagraph:P106:end

\end{minipage}\par\addvspace{4pt}

\par\begin{minipage}{\linewidth}\raggedright
% mdparagraph:P107:begin
\textbf{B4.}
% mdparagraph:P107:end

% item: NEW-U04-B04
% mdparagraph:P108:begin
For sēoðan and drēogan, give the second- and third-person singular present, the first- and second-person singular preterite and the plural preterite. Keep the two verbs’ consonant patterns distinct.
% mdparagraph:P108:end

\end{minipage}\par\addvspace{4pt}

\par\begin{minipage}{\linewidth}\raggedright
% mdparagraph:P109:begin
\textbf{B5.}
% mdparagraph:P109:end

% item: NEW-U04-B05
% mdparagraph:P110:begin
Give the complete selected present and preterite indicative of tēon, “draw”. Use the Class II verb assigned in this unit.
% mdparagraph:P110:end

\end{minipage}\par\addvspace{4pt}

\par\begin{minipage}{\linewidth}\raggedright
% mdparagraph:P111:begin
\textbf{B6.}
% mdparagraph:P111:end

% item: NEW-U04-B06
% mdparagraph:P112:begin
Compare the assigned wrēon and tēon tables. Give each verb’s first-person singular present, first-person singular preterite, second-person singular preterite and plural preterite. Explain why the similar infinitive endings do not make them members of the same strong class.
% mdparagraph:P112:end

\end{minipage}\par\addvspace{4pt}

\reserve{7\baselineskip}
% mdparagraph:P113:begin
\subsection*{Class III — C1–C8}
% mdparagraph:P113:end

\reserve{5\baselineskip}
\par\begin{minipage}{\linewidth}\raggedright
% mdparagraph:P114:begin
\textbf{C1.}
% mdparagraph:P114:end

% item: NEW-U04-C01
% mdparagraph:P115:begin
Give the complete selected present and preterite indicative of bindan. Explain the selected second- and third-person singular present forms.
% mdparagraph:P115:end

\end{minipage}\par\addvspace{4pt}

\par\begin{minipage}{\linewidth}\raggedright
% mdparagraph:P116:begin
\textbf{C2.}
% mdparagraph:P116:end

% item: NEW-U04-C02
% mdparagraph:P117:begin
For helpan and gieldan, give the second- and third-person singular present, the first- and second-person singular preterite and the plural preterite. Preserve each verb’s selected stem vowels.
% mdparagraph:P117:end

\end{minipage}\par\addvspace{4pt}

\par\begin{minipage}{\linewidth}\raggedright
% mdparagraph:P118:begin
\textbf{C3.}
% mdparagraph:P118:end

% item: NEW-U04-C03
% mdparagraph:P119:begin
For weorpan and berstan, give the second- and third-person singular present, the first- and second-person singular preterite and the plural preterite. Explain the difference between their first preterite stems.
% mdparagraph:P119:end

\end{minipage}\par\addvspace{4pt}

\par\begin{minipage}{\linewidth}\raggedright
% mdparagraph:P120:begin
\textbf{C4.}
% mdparagraph:P120:end

% item: NEW-U04-C04
% mdparagraph:P121:begin
Give the complete selected present and preterite indicative of weorðan. Identify the cells with ð and those with d.
% mdparagraph:P121:end

\end{minipage}\par\addvspace{4pt}

\par\begin{minipage}{\linewidth}\raggedright
% mdparagraph:P122:begin
\textbf{C5.}
% mdparagraph:P122:end

% item: NEW-U04-C05
% mdparagraph:P123:begin
Give the complete selected present and preterite indicative of findan. State the ordinary third-person singular preterite selected here and its documented alternative. Then give the second- and third-person singular present and plural preterite of winnan.
% mdparagraph:P123:end

\end{minipage}\par\addvspace{4pt}

\par\begin{minipage}{\linewidth}\raggedright
% mdparagraph:P124:begin
\textbf{C6.}
% mdparagraph:P124:end

% item: NEW-U04-C06
% mdparagraph:P125:begin
For drincan and singan, give the second- and third-person singular present, the first- and second-person singular preterite and the plural preterite. State which Class III pattern they follow.
% mdparagraph:P125:end

\end{minipage}\par\addvspace{4pt}

\par\begin{minipage}{\linewidth}\raggedright
% mdparagraph:P126:begin
\textbf{C7.}
% mdparagraph:P126:end

% item: NEW-U04-C07
% mdparagraph:P127:begin
Give the complete selected present and preterite indicative of tōberstan and forweorðan. Identify the base pattern used by each, and give the separately learned meaning of each whole verb.
% mdparagraph:P127:end

\end{minipage}\par\addvspace{4pt}

\par\begin{minipage}{\linewidth}\raggedright
% mdparagraph:P128:begin
\textbf{C8.}
% mdparagraph:P128:end

% item: NEW-U04-C08
% mdparagraph:P129:begin
Consult the complete reference tables of bregdan and fēolan. Identify their third-person singular present and plural preterite entries and the special changes noted with them. These are recognition reference forms.
% mdparagraph:P129:end

\end{minipage}\par\addvspace{4pt}

\reserve{7\baselineskip}
% mdparagraph:P130:begin
\subsection*{Nouns — N1–N3}
% mdparagraph:P130:end

\reserve{5\baselineskip}
\par\begin{minipage}{\linewidth}\raggedright
% mdparagraph:P131:begin
\textbf{N1.}
% mdparagraph:P131:end

% item: NEW-U04-N01
% mdparagraph:P132:begin
Give the nominative, genitive, dative and accusative singular and plural of scip and wind. Explain their different nominative and accusative plural endings.
% mdparagraph:P132:end

\end{minipage}\par\addvspace{4pt}

\par\begin{minipage}{\linewidth}\raggedright
% mdparagraph:P133:begin
\textbf{N2.}
% mdparagraph:P133:end

% item: NEW-U04-N02
% mdparagraph:P134:begin
Give the nominative, genitive, dative and accusative singular and plural of strand and gefēra. Explain why gefēran alone does not identify a unique case and number.
% mdparagraph:P134:end

\end{minipage}\par\addvspace{4pt}

\par\begin{minipage}{\linewidth}\raggedright
% mdparagraph:P135:begin
\textbf{N3.}
% mdparagraph:P135:end

% item: NEW-U04-N03
% mdparagraph:P136:begin
Translate “The companions found two ships on the shore. The good man’s ship was small.” In the second sentence put the limiting genitive before the head noun.
% mdparagraph:P136:end

\end{minipage}\par\addvspace{4pt}

\reserve{7\baselineskip}
% mdparagraph:P137:begin
\subsection*{Clauses — D1–D10}
% mdparagraph:P137:end

\reserve{5\baselineskip}
\par\begin{minipage}{\linewidth}\raggedright
% mdparagraph:P138:begin
\textbf{D1.}
% mdparagraph:P138:end

% item: NEW-U04-D01
% mdparagraph:P139:begin
Translate “The companions choose the small dwelling. We choose the great ship.”
% mdparagraph:P139:end

\end{minipage}\par\addvspace{4pt}

\par\begin{minipage}{\linewidth}\raggedright
% mdparagraph:P140:begin
\textbf{D2.}
% mdparagraph:P140:end

% item: NEW-U04-D02
% mdparagraph:P141:begin
Translate “The man found the good man’s ship. He bound the animal.” Put the limiting genitive before its head noun.
% mdparagraph:P141:end

\end{minipage}\par\addvspace{4pt}

\par\begin{minipage}{\linewidth}\raggedright
% mdparagraph:P142:begin
\textbf{D3.}
% mdparagraph:P142:end

% item: NEW-U04-D03
% mdparagraph:P143:begin
Translate “We two drank the water, and you two sang well.” Retain both dual subjects.
% mdparagraph:P143:end

\end{minipage}\par\addvspace{4pt}

\par\begin{minipage}{\linewidth}\raggedright
% mdparagraph:P144:begin
\textbf{D4.}
% mdparagraph:P144:end

% item: NEW-U04-D04
% mdparagraph:P145:begin
Change the subjects to the corresponding plurals and make the verbs agree: Se mann bint ðæt dēor. Se fugol singð wel. Keep the object of the first sentence singular.
% mdparagraph:P145:end

\end{minipage}\par\addvspace{4pt}

\par\begin{minipage}{\linewidth}\raggedright
% mdparagraph:P146:begin
\textbf{D5.}
% mdparagraph:P146:end

% item: NEW-U04-D05
% mdparagraph:P147:begin
Change Ic rīde wīde and finde ān hof to the preterite. Preserve the first-person subject, the object and the coordination.
% mdparagraph:P147:end

\end{minipage}\par\addvspace{4pt}

\par\begin{minipage}{\linewidth}\raggedright
% mdparagraph:P148:begin
\textbf{D6.}
% mdparagraph:P148:end

% item: NEW-U04-D06
% mdparagraph:P149:begin
Here hē and the companion are different people. Change Hē sang wel, and se gefēra hīerde hine to the present. Identify the subject of each verb and the referent of hine.
% mdparagraph:P149:end

\end{minipage}\par\addvspace{4pt}

\par\begin{minipage}{\linewidth}\raggedright
% mdparagraph:P150:begin
\textbf{D7.}
% mdparagraph:P150:end

% item: NEW-U04-D07
% mdparagraph:P151:begin
Translate “The companion became glad. The companions become glad.” Explain the case and number of the predicate adjectives.
% mdparagraph:P151:end

\end{minipage}\par\addvspace{4pt}

\par\begin{minipage}{\linewidth}\raggedright
% mdparagraph:P152:begin
\textbf{D8.}
% mdparagraph:P152:end

% item: NEW-U04-D08
% mdparagraph:P153:begin
Translate “I helped the old man. The old man helped me.” Use the case required by the taught helpan construction.
% mdparagraph:P153:end

\end{minipage}\par\addvspace{4pt}

\par\begin{minipage}{\linewidth}\raggedright
% mdparagraph:P154:begin
\textbf{D9.}
% mdparagraph:P154:end

% item: NEW-U04-D09
% mdparagraph:P155:begin
Translate “The companions did not find the ship. I do not choose the great dwelling.”
% mdparagraph:P155:end

\end{minipage}\par\addvspace{4pt}

\par\begin{minipage}{\linewidth}\raggedright
% mdparagraph:P156:begin
\textbf{D10.}
% mdparagraph:P156:end

% item: NEW-U04-D10
% mdparagraph:P157:begin
Translate “You did not drink the water. You drink water.” Address one person in the first sentence and several people in the second.
% mdparagraph:P157:end

\end{minipage}\par\addvspace{4pt}

\reserve{7\baselineskip}
% mdparagraph:P158:begin
\section*{9. Controlled reading: Two companions choose a dwelling}
% mdparagraph:P158:end

% teaching: NEW-U04-GRAM-PAST-NARRATION
% mdparagraph:P159:begin
This account is newly composed. The two companions are the two men of the account; one speaks to the other at the end. Its preterites narrate the successive actions. The quotation shifts to present statements and to the speaker’s first person and addressee’s second person.
% mdparagraph:P159:end

% composition: NEW-U04-TEXT-COMPANIONS-01
% mdparagraph:P160:begin
\par\begin{minipage}{\linewidth}\raggedright
\readingtext{Twēgen gefēran ridon wīde. Hīe fundon twā hofu. Ān hof wæs micel, and ān hof wæs lȳtel. Hīe curon ðæt lȳtle hof. Hīe hæfdon męte and wæter. Hīe druncon ðæt wæter and geręston hīe on ðǣm lȳtlan hofe. Ān mann sægde: “Wit habbað gōd hof. Ic singe, and ðū hīerst mē.”}
\end{minipage}\par
% mdparagraph:P160:end

\reserve{7\baselineskip}
% mdparagraph:P161:begin
\subsection*{Controlled reading — L1–L6}
% mdparagraph:P161:end

\reserve{5\baselineskip}
\par\begin{minipage}{\linewidth}\raggedright
% mdparagraph:P162:begin
\textbf{L1.}
% mdparagraph:P162:end

% item: NEW-U04-L01
% mdparagraph:P163:begin
Translate the complete new account “Two companions choose a dwelling”. Preserve the sequence of actions and the speakers in the closing quotation.
% mdparagraph:P163:end

\end{minipage}\par\addvspace{4pt}

\par\begin{minipage}{\linewidth}\raggedright
% mdparagraph:P164:begin
\textbf{L2.}
% mdparagraph:P164:end

% item: NEW-U04-L02
% mdparagraph:P165:begin
List the strong finite verbs in the account. Give each verb’s headword, class, person, number and tense.
% mdparagraph:P165:end

\end{minipage}\par\addvspace{4pt}

\par\begin{minipage}{\linewidth}\raggedright
% mdparagraph:P166:begin
\textbf{L3.}
% mdparagraph:P166:end

% item: NEW-U04-L03
% mdparagraph:P167:begin
For ridon, fundon, curon and druncon in the account, give the corresponding third-person singular preterite and the second-person singular preterite. Preserve the selected lexical exceptions.
% mdparagraph:P167:end

\end{minipage}\par\addvspace{4pt}

\par\begin{minipage}{\linewidth}\raggedright
% mdparagraph:P168:begin
\textbf{L4.}
% mdparagraph:P168:end

% item: NEW-U04-L04
% mdparagraph:P169:begin
Which dwelling do the companions choose? Who sings in the quotation, and whom does the addressee hear? Explain the reference of wit, ic, ðū and mē.
% mdparagraph:P169:end

\end{minipage}\par\addvspace{4pt}

\par\begin{minipage}{\linewidth}\raggedright
% mdparagraph:P170:begin
\textbf{L5.}
% mdparagraph:P170:end

% item: NEW-U04-L05
% mdparagraph:P171:begin
Retell the first four sentences of the account in the present. Preserve the people, things, case relations and adjective endings.
% mdparagraph:P171:end

\end{minipage}\par\addvspace{4pt}

\par\begin{minipage}{\linewidth}\raggedright
% mdparagraph:P172:begin
\textbf{L6.}
% mdparagraph:P172:end

% item: NEW-U04-L06
% mdparagraph:P173:begin
Replace the closing quotation with “We have a good dwelling. We sing, and you hear us.” Let a group of speakers address another group; use plural personal pronouns.
% mdparagraph:P173:end

\end{minipage}\par\addvspace{4pt}

\reserve{7\baselineskip}
% mdparagraph:P174:begin
\section*{10. Source reading: The shipwreck and the survivor}
% mdparagraph:P174:end

% mdparagraph:P175:begin
This original extract from \textbf{Apollonius of Tyre} gives a complete storm-to-shore episode. The two Class III forms tōbærst and forwurdon use the systems learned here. The additional words and constructions are supplied in the local help and may be consulted throughout these tasks.
% mdparagraph:P175:end

% mdparagraph:P176:begin
\textbf{Setting.} Apollonius has embarked for Cyrene. A sudden storm has already disturbed the sea.
% mdparagraph:P176:end

% source: FB-XII-APOLLONIUS-STORM-SHORE-01
% mdparagraph:P177:begin
\par\begin{minipage}{\linewidth}\raggedright
\readingtext{Ðǣr-tō-ēacan cōmon ēastnorðerne windas, and se angrīslica sūðwesterna wind him ongēan stōd, and ðæt scip eall tōbærst on ðisse ęgeslican hrēohnesse. Apollonius gefēran ealle forwurdon tō dēaðe, and Apollonius āna becōm mid sunde tō Pentapolim ðǣm Cyreniscan lande, and ðǣr ūpēode on ðǣm strande.}
\end{minipage}\par
% mdparagraph:P177:end

\reserve{7\baselineskip}
% mdparagraph:P178:begin
\subsection*{Local help}
% mdparagraph:P178:end

% mdparagraph:P179:begin
The case analyses below apply to the actual phrases. The four supplied preterites cōmon, stōd, becōm and ūpēode do not yet require complete paradigms or independent recall.
% mdparagraph:P179:end

\reserve{5\baselineskip}
% mdparagraph:P180:begin
\textbf{Ðǣr-tō-ēacan}
% mdparagraph:P180:end

% teaching: NEW-U04-SUPFORM-thaerto-eacan-invariant
% teaching: NEW-U04-SUPSENSE-thaerto-eacan-in-addition
% teaching: NEW-U04-GRAM-SOURCE-CONNECTIVE
% teaching: NEW-U04-SUPPORT-U04-SRC-HELP01
% mdparagraph:P181:begin
“In addition to that; besides.” Further winds join the storm already described in the setting.
% mdparagraph:P181:end

\reserve{5\baselineskip}
% mdparagraph:P182:begin
\textbf{cōmon ēastnorðerne windas}
% mdparagraph:P182:end

% teaching: NEW-U04-SUPFORM-cuman-pret-3-pl-source
% teaching: NEW-U04-SUPFORM-eastnortherne-nom-m-pl-source
% teaching: NEW-U04-SUPSENSE-cuman-come
% teaching: NEW-U04-SUPSENSE-eastnortherne-northeasterly
% teaching: NEW-U04-SUPPORT-U04-SRC-HELP02
% mdparagraph:P183:begin
Cōmon is third-person plural preterite, “came”. Ēastnorðerne means “northeasterly” and is masculine nominative plural, agreeing with windas. The winds are the subject.
% mdparagraph:P183:end

\reserve{5\baselineskip}
% mdparagraph:P184:begin
\textbf{se angrīslica sūðwesterna wind him ongēan stōd}
% mdparagraph:P184:end

% teaching: NEW-U04-SUPFORM-angrislic-nom-m-sg-weak-source
% teaching: NEW-U04-SUPFORM-suthwesterne-nom-m-sg-weak-source
% teaching: NEW-U04-SUPFORM-ongean-postposition-invariant-source
% teaching: NEW-U04-SUPFORM-standan-pret-3-sg-source
% teaching: NEW-U04-SUPSENSE-angrislic-fierce
% teaching: NEW-U04-SUPSENSE-suthwesterne-southwestern
% teaching: NEW-U04-SUPSENSE-ongean-against
% teaching: NEW-U04-SUPSENSE-standan-stand
% teaching: NEW-U04-GRAM-SOURCE-OPPOSITION
% teaching: NEW-U04-SUPPORT-U04-SRC-HELP03
% mdparagraph:P185:begin
“The fierce southwest wind stood against them.” Angrīslica “fierce” and sūðwesterna “southwestern” are weak masculine nominative singular after se. Stōd is third-person singular preterite, “stood”. Ongēan follows its dative pronoun him. The guiding interpretation takes him as the preceding northeast winds; its spelling can also be masculine or neuter singular, so the form alone does not settle its reference.
% mdparagraph:P185:end

\reserve{5\baselineskip}
% mdparagraph:P186:begin
\textbf{ðæt scip eall}
% mdparagraph:P186:end

% teaching: NEW-U04-SUPFORM-eall-nom-n-sg-strong
% teaching: NEW-U04-SUPSENSE-eall-all
% teaching: NEW-U04-GRAM-SOURCE-WHOLE-SHIP
% teaching: NEW-U04-SUPPORT-U04-SRC-HELP04
% mdparagraph:P187:begin
“The whole ship; all the ship.” Eall is strong neuter nominative singular and follows the noun phrase. Its morphology remains the recognition material of Unit 3; this whole-ship expression is supplied here.
% mdparagraph:P187:end

\reserve{5\baselineskip}
% mdparagraph:P188:begin
\textbf{on ðisse ęgeslican hrēohnesse}
% mdparagraph:P188:end

% teaching: NEW-U04-SUPFORM-egeslic-dat-f-sg-weak-source
% teaching: NEW-U04-SUPFORM-hreohnes-dat-sg-source
% teaching: NEW-U04-SUPSENSE-on-storm-circumstance
% teaching: NEW-U04-SUPSENSE-egeslic-terrible
% teaching: NEW-U04-SUPSENSE-hreohnes-tempest
% teaching: NEW-U04-GRAM-SOURCE-IN-STORM
% teaching: NEW-U04-SUPPORT-U04-SRC-HELP05
% mdparagraph:P189:begin
“In this terrible tempest.” Hrēohnesse is feminine dative singular; ęgeslican “terrible” is weak after the demonstrative. Here on describes the surrounding circumstance in which the ship breaks up.
% mdparagraph:P189:end

\reserve{5\baselineskip}
% mdparagraph:P190:begin
\textbf{Apollonius gefēran ealle}
% mdparagraph:P190:end

% teaching: NEW-U04-SUPFORM-apollonius-gen-sg-source
% teaching: NEW-U04-SUPFORM-eall-nom-m-pl-strong
% teaching: NEW-U04-SUPSENSE-apollonius-person
% teaching: NEW-U04-GRAM-SOURCE-PROPER-GENITIVE
% teaching: NEW-U04-GRAM-SOURCE-ALL-COMPANIONS
% teaching: NEW-U04-SUPPORT-U04-SRC-HELP06
% mdparagraph:P191:begin
“All Apollonius’s companions.” The first Apollonius is genitive, despite its ending; Cook explicitly identifies its case. Gefēran is nominative plural. Ealle “all” is strong masculine nominative plural and follows the noun. The all-quantifier phrase is supplied, while the noun and its plural agreement use the paradigms already available.
% mdparagraph:P191:end

\reserve{5\baselineskip}
% mdparagraph:P192:begin
\textbf{forwurdon tō dēaðe}
% mdparagraph:P192:end

% teaching: NEW-U04-SUPFORM-to-preposition-dative-invariant
% teaching: NEW-U04-SUPFORM-death-dat-sg-source
% teaching: NEW-U04-SUPSENSE-to-death-result
% teaching: NEW-U04-SUPSENSE-death-death
% teaching: NEW-U04-GRAM-SOURCE-DEATH-ENDPOINT
% teaching: NEW-U04-SUPPORT-U04-SRC-HELP07
% mdparagraph:P193:begin
“Perished to death”; naturally, “died” or “perished”. Forwurdon is the taught third-person plural preterite. Dēaðe “death” is masculine dative singular. The whole tō phrase gives the endpoint or result and is supplied here.
% mdparagraph:P193:end

\reserve{5\baselineskip}
% mdparagraph:P194:begin
\textbf{Apollonius āna becōm mid sunde}
% mdparagraph:P194:end

% teaching: NEW-U04-SUPFORM-apollonius-nom-sg-source
% teaching: NEW-U04-SUPFORM-an-nom-m-sg-weak-alone-source
% teaching: NEW-U04-SUPFORM-becuman-pret-3-sg-source
% teaching: NEW-U04-SUPFORM-mid-invariant-source
% teaching: NEW-U04-SUPFORM-sund-dat-sg-source
% teaching: NEW-U04-SUPSENSE-an-alone
% teaching: NEW-U04-SUPSENSE-becuman-arrive
% teaching: NEW-U04-SUPSENSE-mid-by-means
% teaching: NEW-U04-SUPSENSE-sund-swimming
% teaching: NEW-U04-GRAM-SOURCE-ALONE
% teaching: NEW-U04-GRAM-SOURCE-BY-SWIMMING
% teaching: NEW-U04-SUPPORT-U04-SRC-HELP08
% mdparagraph:P195:begin
“Apollonius alone arrived by swimming.” This Apollonius is nominative. Āna is weak masculine nominative singular in the sense “alone”, distinct from the ordinary cardinal “one”. Becōm is third-person singular preterite, “came, arrived”. Sunde is the dative singular of a neuter noun meaning “swimming”; mid sunde gives the means. It is not an infinitive or an instrumental case form.
% mdparagraph:P195:end

\reserve{5\baselineskip}
% mdparagraph:P196:begin
\textbf{tō Pentapolim ðǣm Cyreniscan lande}
% mdparagraph:P196:end

% teaching: NEW-U04-SUPFORM-pentapolis-goal-source
% teaching: NEW-U04-SUPFORM-cyrenisc-dat-n-sg-weak-source
% teaching: NEW-U04-SUPFORM-land-dat-sg-source
% teaching: NEW-U04-SUPSENSE-to-destination
% teaching: NEW-U04-SUPSENSE-pentapolis-place
% teaching: NEW-U04-SUPSENSE-cyrenisc-of-cyrene
% teaching: NEW-U04-SUPSENSE-land-country
% teaching: NEW-U04-GRAM-SOURCE-DESTINATION
% teaching: NEW-U04-SUPPORT-U04-SRC-HELP09
% mdparagraph:P197:begin
“To Pentapolis, the Cyrenian land.” The foreign place-name spelling Pentapolim is preserved. The native phrase that explains the destination is neuter dative singular; Cyreniscan is weak and means “Cyrenian, of Cyrene”; lande means “land, country”. This exact destination and apposition are supplied without a new general rule for foreign-name endings.
% mdparagraph:P197:end

\reserve{5\baselineskip}
% mdparagraph:P198:begin
\textbf{ðǣr ūpēode on ðǣm strande}
% mdparagraph:P198:end

% teaching: NEW-U04-SUPFORM-thaer-invariant-source
% teaching: NEW-U04-SUPFORM-upgan-pret-3-sg-source
% teaching: NEW-U04-SUPSENSE-thaer-there
% teaching: NEW-U04-SUPSENSE-upgan-go-up
% teaching: NEW-U04-GRAM-SOURCE-SHORE-ARRIVAL
% teaching: NEW-U04-SUPPORT-U04-SRC-HELP10
% mdparagraph:P199:begin
“There he came up onto the shore.” Ðǣr means “there”; ūpēode is the actual third-person singular preterite, “went/\allowbreak{}came up”. Apollonius remains the understood subject. Strande is dative singular in the printed expression, even though a natural English translation describes movement.
% mdparagraph:P199:end

\reserve{7\baselineskip}
% mdparagraph:P200:begin
\subsection*{Source reading — R1–R5}
% mdparagraph:P200:end

\reserve{5\baselineskip}
\par\begin{minipage}{\linewidth}\raggedright
% mdparagraph:P201:begin
\textbf{R1.}
% mdparagraph:P201:end

% item: NEW-U04-R01
% mdparagraph:P202:begin
Translate the complete source passage “The shipwreck and the survivor”, using the local support. Preserve the contrasting fates of the ship, the companions and Apollonius.
% mdparagraph:P202:end

\end{minipage}\par\addvspace{4pt}

\par\begin{minipage}{\linewidth}\raggedright
% mdparagraph:P203:begin
\textbf{R2.}
% mdparagraph:P203:end

% item: NEW-U04-R02
% mdparagraph:P204:begin
Analyse tōbærst and forwurdon: give each whole verb’s headword, class, person, number and tense, and identify its subject in the passage.
% mdparagraph:P204:end

\end{minipage}\par\addvspace{4pt}

\par\begin{minipage}{\linewidth}\raggedright
% mdparagraph:P205:begin
\textbf{R3.}
% mdparagraph:P205:end

% item: NEW-U04-R03
% mdparagraph:P206:begin
Using the proper-name note, distinguish the first Apollonius from the second. Explain why gefēran takes plural forwurdon rather than a singular verb.
% mdparagraph:P206:end

\end{minipage}\par\addvspace{4pt}

\par\begin{minipage}{\linewidth}\raggedright
% mdparagraph:P207:begin
\textbf{R4.}
% mdparagraph:P207:end

% item: NEW-U04-R04
% mdparagraph:P208:begin
Identify windas, wind, scip, gefēran and strande in their clauses. Give each noun’s case and number; use the supplied phrase notes wherever the construction is new.
% mdparagraph:P208:end

\end{minipage}\par\addvspace{4pt}

\par\begin{minipage}{\linewidth}\raggedright
% mdparagraph:P209:begin
\textbf{R5.}
% mdparagraph:P209:end

% item: NEW-U04-R05
% mdparagraph:P210:begin
How does Apollonius reach land, and where does the passage leave him? Explain the contribution of the supplied āna without treating it as an ordinary cardinal “one”.
% mdparagraph:P210:end

\end{minipage}\par\addvspace{4pt}

\reserve{7\baselineskip}
% mdparagraph:P211:begin
\section*{11. Cumulative review}
% mdparagraph:P211:end

\reserve{7\baselineskip}
% mdparagraph:P212:begin
\subsection*{Review — X1–X4}
% mdparagraph:P212:end

\reserve{5\baselineskip}
\par\begin{minipage}{\linewidth}\raggedright
% mdparagraph:P213:begin
\textbf{X1.}
% mdparagraph:P213:end

% item: NEW-U04-X01
% mdparagraph:P214:begin
Correct the three marked finite forms: Ðū *drāf. Hīe *cēas ðæt scip. Se mann *bindað ðæt dēor. Preserve the given subjects and tenses.
% mdparagraph:P214:end

\end{minipage}\par\addvspace{4pt}

\par\begin{minipage}{\linewidth}\raggedright
% mdparagraph:P215:begin
\textbf{X2.}
% mdparagraph:P215:end

% item: NEW-U04-X02
% mdparagraph:P216:begin
Write four connected Old English sentences about two companions: they ride, find two ships, choose the small ship, and drink water. Narrate all four actions in the preterite and retain the same two people as subjects.
% mdparagraph:P216:end

\end{minipage}\par\addvspace{4pt}

\par\begin{minipage}{\linewidth}\raggedright
% mdparagraph:P217:begin
\textbf{X3.}
% mdparagraph:P217:end

% item: NEW-U04-X03
% mdparagraph:P218:begin
Translate “The old companion helps my birds. My birds sing well. We found his small ship on the shore.” Distinguish the helper, the birds, the finders and the owner.
% mdparagraph:P218:end

\end{minipage}\par\addvspace{4pt}

\par\begin{minipage}{\linewidth}\raggedright
% mdparagraph:P219:begin
\textbf{X4.}
% mdparagraph:P219:end

% item: NEW-U04-X04
% mdparagraph:P220:begin
Analyse Hē funde, Ðū funde and Hīe fundon. Which finite forms distinguish number? What identifies the different persons in the two singular clauses?
% mdparagraph:P220:end

\end{minipage}\par\addvspace{4pt}

\clearpage
% mdparagraph:P221:begin
\section*{12. Source and variant notes}
% mdparagraph:P221:end

\begingroup\fontsize{10}{12.8}\selectfont
\setlength{\parskip}{1.5pt}
% mdparagraph:P222:begin
\itempara{1}{The principal-part, indicative and stem rules follow Cook, \textbf{A First Book in Old English}, §§92–95 and 101–104, printed pp.54–55 and 58–63. The selected tables preserve the edition’s quantities. Full finite cells formed from a checked stem and the ordinary endings are transparent applications of the pattern; the tables do not claim a separate surviving quotation for every person.}
% mdparagraph:P222:end

% mdparagraph:P223:begin
\itempara{2}{Wright’s \textbf{Old English Grammar} corroborates the selected stem series and recorded alternatives. In particular, its Class III discussion permits first- as well as third-person funde. Its discussion of contracted Class I verbs records analogical movement into Class II; the wrēon table here consistently selects Cook’s Class I pattern. Alternatives in the notes are available for recognition and comparison, not additional independent-recall requirements.}
% mdparagraph:P223:end

% mdparagraph:P224:begin
\itempara{3}{The rare fēolan reference combines Cook’s principal parts and ordinary endings with Bosworth–Toller’s filhst/\allowbreak{}filhþ entries. The displayed filhð changes þ to the course’s ð spelling only. The supplement transfers the earlier entry’s passages to fēolan. This is a complete synthesized reference, not a claim that each printed cell is individually attested. Cook’s folen and Wright’s folgen/\allowbreak{}fōlen remain distinct citations.}
% mdparagraph:P224:end

% mdparagraph:P225:begin
\itempara{4}{The ordinary case frames are independently documented: Cook’s direct-object and predicate rules (§§168 and 150), his occurrences of findan, cēosan and singan, and the Bosworth–Toller entries for drīfan, bindan, rīdan, drincan and helpan. The animal and ship/\allowbreak{}dwelling objects in the new examples are ordinary applications of these senses. Attested human dative recipients of help support its ordinary application to birds; the dictionary also records a genitive construction.}
% mdparagraph:P225:end

% mdparagraph:P226:begin
\itempara{5}{The original extract is Cook, \textbf{A First Book in Old English}, Reader XII, \textbf{Apollonius of Tyre}, printed pp.165–166, one-based PDF181–182. It begins within p.165 line 6 and ends at the first period in p.166 line 4. These line counts expand the edition’s printed five-line intervals. The line-crossing compound Ðǣr-tō-ēacan is joined, and note cues are moved below; the Old English spelling and punctuation are preserved. The English setting is editorial context.}
% mdparagraph:P226:end

\reserve{7\baselineskip}
% mdparagraph:P227:begin
\subsection*{Cook’s notes to the extract}
% mdparagraph:P227:end

% source-note: FB165-n7
% mdparagraph:P228:begin
\textbf{P.165, note 7.} Lat. (verse): Hinc Notus,hinc Boreas,hinc horridus Africus instat.
% mdparagraph:P228:end

% mdparagraph:P229:begin
The Latin verse evokes winds from different quarters; it does not explicitly settle the antecedent of him.
% mdparagraph:P229:end

% source-note: FB166-n1
% mdparagraph:P230:begin
\textbf{P.166, note 1.} The Latin endings of proper nouns are not always a guide to the case (54). Here we have the genitive.
% mdparagraph:P230:end

% source-note: FB166-n2
% mdparagraph:P231:begin
\textbf{P.166, note 2.} See forweorðan.
% mdparagraph:P231:end

% source-note: FB166-n3
% mdparagraph:P232:begin
\textbf{P.166, note 3.} See 79.
% mdparagraph:P232:end

% mdparagraph:P233:begin
Section 79 explains the special “alone” use supplied with āna above.
% mdparagraph:P233:end

\endgroup
\end{document}
